% ── SECTION 5: NUMERICAL SIMULATION RESULTS (~1,500 words) ──────────────────

\section{Numerical Simulation}
\label{sec:sim}

\subsection{Implementation}

We implement $\URA(\theta)$ as a two-qubit unitary operator via the
explicit matrix form of Eq.~\eqref{eq:URA-matrix} in both
Qiskit~\cite{Qiskit2024} (IBM) and Cirq~\cite{Cirq2025} (Google).
Simulation parameters follow Appendix~\ref{app:sim-params}.
The circuit representation is:

\begin{center}
\begin{quantikz}
  \lstick{$\ket{0}$} & \gate[2]{U_{\mathrm{RA}}(\theta)} & \qw \\
  \lstick{$\ket{0}$} &                                    & \qw
\end{quantikz}
\end{center}

\noindent
Performance is benchmarked against three fixed two-qubit gates~---~\textsc{cnot},
\textsc{iswap}, and \textsc{cz}~---~which are the standard entangling
primitives on IBM and Google hardware respectively.
The key distinction is that these fixed gates always operate at full
entanglement capability ($\mathcal{E} = 1$), whereas $\URA(\theta)$
operates at $\mathcal{E}(\theta) = |\sin 2\theta|$
(Eq.~\eqref{eq:entcap-result}), continuously tunable via the exchange
angle.

\subsection{Noise Models and Gate-Time Scaling}

Three noise channels model the dominant decoherence mechanisms in
near-term superconducting qubit devices:

\begin{enumerate}
  \item \textbf{Depolarizing noise:}
    $\mathcal{E}(\rho) = (1-p)\rho + \tfrac{p}{4}\mathbf{I}_4$,
    where $p \in [0.01, 0.10]$ is the per-gate error rate.
  \item \textbf{Dephasing (phase damping):}
    $\rho_{ij} \mapsto \rho_{ij}(1-\gamma)^{d_{ij}}$,
    where $d_{ij}$ is the Hamming distance between bitstrings $i$ and $j$
    and $\gamma \in [0.02, 0.20]$.
  \item \textbf{Amplitude damping:}
    Kraus channel with $T_1$ relaxation probability
    $p_{\mathrm{AD}} \in [0.02, 0.20]$, applied independently to each qubit.
\end{enumerate}

A critical physical consideration governs the comparison.
On superconducting qubit hardware, two-qubit gate error accumulates in
proportion to gate duration~\cite{Krantz2019}.
The Heisenberg exchange interaction implementing $\URA(\theta)$ has
duration proportional to~$\theta$: running at $\theta = \pi/8$ takes
half the time of a full $\theta = \pi/4$ (\textsc{iswap}) operation.
We therefore apply \emph{gate-time-proportional noise}: fixed gates
(\textsc{cnot}, \textsc{iswap}, \textsc{cz}) receive error rate~$p_{\rm base}$
at full duration, while $\URA(\theta)$ receives effective rate
\begin{equation}
  p_{\rm eff}(\theta) = p_{\rm base} \times \frac{\theta}{\pi/4},
  \label{eq:peff}
\end{equation}
reflecting shorter physical operation time for $\theta < \pi/4$.

\subsection{Results}

\textbf{Gate fidelity advantage.}
Table~\ref{tab:fidelity} reports the average gate fidelity
\begin{equation}
  F_{\rm avg}(W, V) \;=\; \frac{d\,F_{\rm proc}(W, V) + 1}{d + 1},
  \qquad
  F_{\rm proc}(W, V) \;=\; \mathrm{Tr}[\chi_V^\dagger\,\chi_W],
  \label{eq:favg}
\end{equation}
of each gate under depolarizing noise with gate-time-proportional
error scaling, where $d = 4$ for two qubits and $\chi_V$, $\chi_W$ are
the Choi--Jamio\l{}kowski process matrices of the ideal and noisy
implementations respectively.
For the global depolarizing channel $W(\rho) = (1-p_{\rm eff})\rho +
p_{\rm eff}\,\mathbf{I}_4/4$ applied after an ideal unitary, this
reduces to the closed-form expression
$F_{\rm avg} = 1 - 3\,p_{\rm eff}/4$, which we report directly;
\texttt{figures/generate\_tables.py} provides an analytical
sanity-check, and the Qiskit/Cirq density-matrix simulator notebooks
of the release tag (Sec.~\ref{sec:conc}) confirm the values to numerical
precision.
$\URA(\pi/4)$ (full duration, $\theta = \pi/4$) operates at maximum
entanglement capability ($\mathcal{E} = 1$) and, as the maximally entangling
point of the Heisenberg diagonal (Sec.~\ref{subsec:weyl}), matches the
fidelity of \textsc{iswap} under gate-time-proportional noise at equal duration.
$\URA(\pi/8)$ (half duration, $\theta = \pi/8$) achieves measurable
advantages: $+1.87\%$ at $p_{\rm base} = 0.05$ and $+3.75\%$ at
$p_{\rm base} = 0.10$ over all three fixed gates under depolarizing noise.

\begin{table}[htbp]
\centering
\caption{Average gate fidelity under depolarizing noise (gate-time-proportional
model). $p_{\rm base}$ = error rate per full-duration two-qubit gate.
$\URA(\pi/8)$ operates at $0.50\times$ gate duration; $\URA(\pi/6)$ at
$0.67\times$.}
\label{tab:fidelity}
\begin{tabular}{lccccl}
\toprule
Gate & $p=0.01$ & $p=0.02$ & $p=0.05$ & $p=0.10$ & Duration \\
\midrule
\textsc{cnot}    & 0.9925 & 0.9850 & 0.9625 & 0.9250 & $1.00\times$ \\
\textsc{iswap}   & 0.9925 & 0.9850 & 0.9625 & 0.9250 & $1.00\times$ \\
\textsc{cz}      & 0.9925 & 0.9850 & 0.9625 & 0.9250 & $1.00\times$ \\
\midrule
$\URA(\pi/4)$    & 0.9925 & 0.9850 & 0.9625 & 0.9250 & $1.00\times$ \\
$\URA(\pi/6)$    & 0.9950 & 0.9900 & 0.9750 & 0.9500 & $0.67\times$ \\
$\URA(\pi/8)$    & \textbf{0.9962} & \textbf{0.9925} & \textbf{0.9812} & \textbf{0.9625} & $0.50\times$ \\
\bottomrule
\end{tabular}
\end{table}

\textbf{Consistency across noise models.}
Table~\ref{tab:fidelity-all} reports the fidelity advantage
$\Delta F = F[\URA(\pi/8)] - F[\textsc{iswap}]$ at $p_{\rm base} \in \{0.05, 0.10\}$
across all three noise models.
The advantage is consistent: $\Delta F \in [+0.019, +0.040]$ at
$p_{\rm base} = 0.10$ across depolarizing, dephasing, and amplitude damping.

\begin{table}[htbp]
\centering
\caption{Fidelity advantage $\Delta F = F[\URA(\pi/8)] - F[\mathrm{iSWAP}]$
under gate-time-proportional noise. Positive values confirm that
$\URA(\pi/8)$ outperforms the fixed gates in all high-decoherence regimes tested.}
\label{tab:fidelity-all}
\begin{tabular}{lccc}
\toprule
Noise model & $p_{\rm base} = 0.02$ & $p_{\rm base} = 0.05$ & $p_{\rm base} = 0.10$ \\
\midrule
Depolarizing    & $+0.0075$ & $+0.0187$ & $+0.0375$ \\
Dephasing       & $+0.0079$ & $+0.0197$ & $+0.0389$ \\
Amplitude damp. & $+0.0080$ & $+0.0195$ & $+0.0395$ \\
\bottomrule
\end{tabular}
\end{table}

\textbf{Interpretation via the \RCT.}
The fidelity advantage is naturally interpreted through the
conservation law of Proposition~\ref{prop:conservation}.
A gate operating at $\theta < \pi/4$ implements a weaker exchange
interaction and consequently requires less physical interaction time
with the bath, accumulating less decoherence per operation.
Operationally, this means a smaller fraction of the local pointer-basis
coherence budget $\bar\alpha$ (Definition~\ref{def:alpha-scalar}) is
expended into $S$-bath correlations during the gate, leaving more
coherence available for downstream layers~---~exactly the redistribution
quantified by Eq.~\eqref{eq:coherence-transfer}.
For circuits requiring partial entanglement, or in which entanglement
is built incrementally across layers, $\URA(\theta)$ can therefore be
tuned to the coherence regime of the device: small $\theta$ in
high-noise environments preserves coherence that a fixed maximally
entangling gate would dissipate.
The \RCT~contributes the conceptual warrant for this design~---~that
relational coupling is a continuous, conservatively redistributed
resource rather than a binary primitive~---~while the numerics confirm
the predicted gate-time scaling.

\textbf{$\bar\alpha(t)$ dynamics.}
Figure~\ref{fig:alpha-dynamics} shows the decay of $\bar\alpha(t)$ for a
two-qubit system $S$ under independent per-qubit dephasing at
$\gamma = 0.10$, initialized in the equal-weight pointer-basis
superposition
$\ket{++} = \tfrac{1}{2}(\ket{00}+\ket{01}+\ket{10}+\ket{11})$, for
which $\bar\alpha(0) = 1$ in the four-dimensional computational pointer
basis (Proposition~\ref{prop:correspondence}\emph{(ii)}).
The Affinity Tensor decays monotonically from $1$ to $0$, tracking the
loss of local pointer-basis coherence as the bath selects the classical
pointer basis. Per
Proposition~\ref{prop:conservation}, purity transfers from the system
to $S$-bath correlations: although $\rho_S$ becomes mixed, the joint
$\mathrm{Tr}[\rho_{S,\mathrm{bath}}^2]$ remains constant in the closed
$S$+bath description.

\begin{figure}[htbp]
  \centering
  \includegraphics[width=0.85\columnwidth]{alpha-dynamics}
  \caption{Affinity Tensor dynamics $\bar\alpha(t)$ for the two-qubit
    initial state $\ket{++}$ under independent per-qubit dephasing
    ($\gamma = 0.10$). The monotonic decay from $\bar\alpha = 1$ (maximal
    local pointer-basis coherence) to $\bar\alpha = 0$ (pointer-basis
    diagonal $\rho_S$) tracks the quantum-to-classical transition for
    the local state. The closed-form curve
    $\bar\alpha(t) = \sqrt{\tfrac{2}{3}e^{-2\gamma t}
                           + \tfrac{1}{3}e^{-4\gamma t}}$ reflects the
    two off-diagonal classes (Hamming distance $1$ and $2$) decaying at
    distinct rates; the dashed reference $e^{-\gamma t}$ is the
    single-qubit decay law of Fig.~\ref{fig:affinity-decay}. Total
    purity of the closed $S$+bath system is conserved throughout
    (Proposition~\ref{prop:conservation}); coherence is redistributed
    into $S$-bath correlations rather than annihilated. $\bar\alpha$
    measures local pointer-basis coherence in the reduced state of the
    full two-qubit system $\rho_S$, not $S$-bath entanglement;
    initializing instead in a Bell state $\ket{\Phi^+}_{q_1q_2}$
    (single-qubit marginals $= \mathbf{I}/2$, intra-$S$ entanglement
    maximal) yields $\bar\alpha(0) = \sqrt{2/3}$ in the four-dimensional
    computational pointer basis, not $1$.
    Generated by \texttt{figures/generate\_fig2.py}.}
  \label{fig:alpha-dynamics}
\end{figure}
